\section{Bài toán nghiệp vụ và dữ liệu}
\subsection{Tổng quan giải pháp}
Để giải quyết vấn đề kinh doanh ở mục 2.3, nhóm đề xuất xây dựng một hệ thống hỗ trợ ra quyết định gồm bốn thành phần nối tiếp nhau:

\begin{itemize}
  \item \textbf{Tích hợp dữ liệu:} trích xuất dữ liệu sự kiện từ nguồn và đưa về một cấu trúc thống nhất.
  \item \textbf{Kho dữ liệu:} lưu dữ liệu đã chuẩn hóa theo mô hình chiều phục vụ phân tích theo thời gian, kênh truy cập, chiến dịch, thiết bị, khu vực địa lý và sản phẩm.
  \item \textbf{Hệ thống chỉ số (KPI):} tính toán các chỉ số theo một bộ định nghĩa duy nhất, dùng chung cho mọi báo cáo.
  \item \textbf{Ứng dụng hỗ trợ quyết định:} dashboard có bộ lọc, phân tích chi tiết, so sánh hiệu suất, cảnh báo theo ngưỡng và mô phỏng kịch bản.
\end{itemize}

Giải pháp hướng đến ba kết quả: (i) mọi con số trong báo cáo đều có thể truy vết được về dữ liệu gốc; (ii) các kênh và chiến dịch được so sánh trên cùng một định nghĩa chỉ số; (iii) người quản lý chuyển được từ quan sát sang hành động, chẳng hạn phát hiện kênh có chất lượng truy cập suy giảm, xác định bước khách hàng rời bỏ nhiều nhất trong phễu mua hàng, hoặc xác định kênh nào đáng ưu tiên đầu tư nội dung và nguồn lực. \medskip

Do dữ liệu sử dụng không chứa chi phí quảng cáo, giải pháp đánh giá hiệu quả dựa trên chất lượng lưu lượng và kết quả kinh doanh (chuyển đổi, doanh thu, hành vi) thay vì chi phí.

\subsection{Tập dữ liệu sử dụng}
\subsubsection{Nguồn dữ liệu}

Dự án sử dụng bộ dữ liệu công khai \texttt{ga4\_obfuscated\_sample\_ecommerce} thuộc chương trình BigQuery Public Datasets của Google. Bộ dữ liệu chứa mẫu dữ liệu sự kiện (event export) đã được làm mờ của Google Merchandise Store, một cửa hàng trực tuyến bán sản phẩm mang thương hiệu Google. Dữ liệu trải dài ba tháng, từ 01/11/2020 đến 31/01/2021. Website triển khai Google Analytics 4 theo chuẩn thương mại điện tử cùng cơ chế đo lường nâng cao (enhanced measurement). Dữ liệu được lưu trong các bảng sự kiện theo ngày \parencite{ga4-export-schema}.

\subsubsection{Cấu trúc và nội dung}

Mỗi dòng là một sự kiện của người dùng. Các nhóm thông tin chính gồm:

\begin{table}[htbp]
    \centering
    \renewcommand{\arraystretch}{1.3}
    \caption{Cấu trúc và nội dung chính của tập dữ liệu}
    \label{tab:dataset_structure}
    \begin{tabular}{|p{2.8cm}|p{6.2cm}|p{5cm}|}
        \hline
        \textbf{Nhóm thông tin} &
        \textbf{Trường tiêu biểu} &
        \textbf{Mục đích phân tích} \\
        \hline
        Thời gian &
        \texttt{event\_date}, \texttt{event\_timestamp} &
        Xu hướng theo ngày, tuần, tháng \\
        \hline
        Người dùng và phiên &
        \texttt{user\_pseudo\_id}, \texttt{ga\_session\_id} (trong \texttt{event\_params}) &
        Đếm người dùng, phiên; phân tích hành vi theo phiên \\
        \hline
        Sự kiện &
        \texttt{event\_name} (\texttt{page\_view}, \texttt{view\_item}, \texttt{add\_to\_cart}, \texttt{begin\_checkout}, \texttt{purchase}, ...) &
        Xây dựng phễu chuyển đổi \\
        \hline
        Nguồn truy cập &
        \texttt{traffic\_source.source}, \texttt{medium}, \texttt{name} &
        So sánh kênh và chiến dịch \\
        \hline
        Thiết bị, địa lý &
        \texttt{device}, \texttt{geo} &
        Phân khúc khách hàng \\
        \hline
        Thương mại điện tử &
        \texttt{ecommerce} (giá trị giao dịch, mã giao dịch), \texttt{items} &
        Doanh thu, đơn hàng, sản phẩm \\
        \hline
    \end{tabular}
\end{table}

\subsubsection{Đặc điểm và giới hạn}

\begin{itemize}
    \item \textbf{Dữ liệu đã bị làm mờ.} Một số trường chứa giá trị giữ chỗ như \texttt{<Other>}, \texttt{NULL} hoặc chuỗi rỗng, và tính nhất quán nội bộ có thể bị hạn chế. Đây là lý do quy trình ETL cần có bước kiểm tra, xử lý và ghi nhận chất lượng dữ liệu.

    \item \textbf{Kết quả khảo sát chất lượng dữ liệu.} Khảo sát toàn bộ 92 ngày (22/10/2026, bộ truy vấn \texttt{survey/01}--\texttt{09}) ghi nhận: 100\% sự kiện có tham số \texttt{ga\_session\_id}; 845 phiên (0,23\%) phát sinh sự kiện trên hai ngày lịch, nên phiên được gán theo ngày sự kiện đầu tiên; 84,36\% người dùng có đúng một cặp source/medium trong toàn kỳ và 15,64\% đổi nguồn từ hai đến tám lần, nên nguồn cấp người dùng được khóa theo sự kiện đầu tiên; khoảng 41\% người dùng có nguồn ẩn danh do làm mờ dữ liệu. Trong 5.692 sự kiện \texttt{purchase}, 23 sự kiện thiếu \texttt{transaction\_id}, chỉ có 4.452 mã giao dịch phân biệt do trùng lặp, và \texttt{add\_to\_cart} vắng mặt trong 18/92 ngày. Các con số này là cơ sở cho quy tắc giao dịch hợp lệ ở mục 3.4 và cho nhật ký chất lượng dữ liệu.

    \item \textbf{Không có chi phí quảng cáo.} Bộ dữ liệu không chứa chi phí, lượt hiển thị hay lượt nhấp từ nền tảng quảng cáo, nên các chỉ số như CPA, ROAS, CPC không nằm trong phạm vi đồ án. Nhóm không tự tạo dữ liệu chi phí giả định, mà đánh giá hiệu quả kênh và chiến dịch bằng các chỉ số \textbf{không phụ thuộc chi phí}: lưu lượng, hành vi, tỷ lệ chuyển đổi, doanh thu và giá trị đơn hàng. Khi doanh nghiệp có dữ liệu chi phí thực, bảng chi phí có thể được bổ sung vào kho theo cùng chiều thời gian và chiều kênh/chiến dịch để mở rộng sang CPA, ROAS mà không phải thay đổi kiến trúc.

    \item \textbf{Phạm vi thời gian ngắn.} Dữ liệu chỉ trải dài ba tháng, bao gồm mùa mua sắm cuối năm, nên các kết luận về xu hướng và tính mùa vụ cần được diễn giải thận trọng.
\end{itemize}

Tập dữ liệu được lựa chọn vì mô phỏng sát một doanh nghiệp thương mại điện tử có website, nhiều nguồn truy cập và luồng mua hàng đầy đủ, phù hợp với bối cảnh và đáp ứng yêu cầu dùng bộ dữ liệu công khai. Phạm vi dữ liệu đầu vào này không giả định thêm chi phí quảng cáo. Các hướng mở rộng dữ liệu và hạ tầng được trình bày riêng ở phần kết luận.

\subsection{Kiến trúc hệ thống}

Kiến trúc gồm năm tầng chính được tổ chức theo từng chức năng nhằm giảm sự phụ thuộc giữa quá trình thu thập dữ liệu, xử lý nghiệp vụ và khai thác dữ liệu:

\begin{enumerate}
    \item \textbf{Tầng nguồn dữ liệu:} Các bảng sự kiện theo ngày có tên dạng \texttt{events\_YYYYMMDD}. Mỗi bản ghi đại diện cho một sự kiện trên website như xem trang, xem sản phẩm, thêm sản phẩm vào giỏ hàng hoặc mua hàng. Hệ thống sử dụng cơ chế phát lại từng sự kiện để mô phỏng quá trình tiếp nhận dữ liệu trực tuyến.

    \item \textbf{Tầng thu thập (Extract):} Python Publisher truy vấn dữ liệu từ BigQuery theo khoảng thời gian cấu hình, chuyển từng bản ghi thành thông điệp JSON và gửi đến Pub/Sub Emulator trên máy phát triển \parencite{pubsub-emulator}. Hàng đợi này tách quá trình phát lại khỏi quá trình làm sạch dữ liệu, nhưng chỉ được dùng cho thử nghiệm trong phiên emulator, không phải kho lưu trữ dài hạn. Các thông tin khóa và dấu thời gian được giữ để truy vết và kiểm tra trùng lặp. Thông điệp không xử lý được được gửi tới dead-letter topic kèm lý do.

    \item \textbf{Tầng xử lý (Transform):} Apache Beam chạy pipeline bằng DirectRunner trên máy phát triển \parencite{apache-direct-runner}. Pipeline đọc subscription của emulator, biến đổi dữ liệu JSON thành sự kiện hợp lệ rồi ghi các nhánh kết quả ra file JSONL cục bộ. DirectRunner được dùng để kiểm chứng tính đúng đắn của pipeline, không phải để thay thế một hệ thống xử lý phân tán quy mô lớn.
    \begin{itemize}
        \item Giải mã thông điệp JSON và kiểm tra các trường bắt buộc (ngày, tên, thời điểm sự kiện); thông điệp không giải mã được hoặc sai cấu trúc được ghi vào bảng dead-letter kèm lý do, để tránh làm gián đoạn toàn bộ pipeline.
        \item Ép kiểu chuẩn cho dấu thời gian và các trường số, thay giá trị giữ chỗ (\texttt{<Other>}, \texttt{(data deleted)}, \texttt{(not set)}, chuỗi rỗng) bằng giá trị thiếu, và điền cấu trúc rỗng cho \texttt{event\_params}, \texttt{items}, \texttt{ecommerce} để mọi bản ghi có cùng hình dạng. Bản ghi được giữ nguyên cấu trúc lồng nhau của dữ liệu gốc.
        \item Hỗ trợ phát hiện các sự kiện trùng lặp dựa theo bộ trường định danh trong tệp cấu hình chất lượng (\texttt{user\_pseudo\_id}, \texttt{event\_timestamp}, \texttt{event\_name}, \texttt{event\_bundle\_sequence\_id}) trong từng cửa sổ thời gian; bản trùng bị loại và số dòng bị loại được ghi vào nhật ký chất lượng dữ liệu. Việc khử trùng toàn cục tiếp tục được bảo đảm ở tầng staging.
        \item Áp dụng bộ quy tắc chất lượng đồng nhất với tầng kho dữ liệu: loại bỏ và ghi nhật ký đối với các sự kiện \texttt{purchase} thiếu \texttt{transaction\_id} hoặc có doanh thu không hợp lệ ($\le 0$); riêng các sự kiện thiếu định danh phiên vẫn được giữ lại để ánh xạ về bản ghi mặc định ``\texttt{Unknown}'' ở các bảng dimension tầng sau.
        \item Xuất các bản ghi hợp lệ, dead-letter và nhật ký chất lượng ra các tệp JSONL. Sau đó, script \texttt{load\_raw\_to\_bigquery.py} gom nhóm sự kiện theo ngày và nạp vào bảng \texttt{events\_YYYYMMDD} trên BigQuery Sandbox thông qua client-side load job, giúp bảo toàn cấu trúc dữ liệu lồng nhau phục vụ cho tầng staging.
        \item Việc làm phẳng \texttt{event\_params}, tách \texttt{items} thành từng dòng mặt hàng và ghép \texttt{user\_pseudo\_id} với \texttt{ga\_session\_id} thành định danh phiên được thực hiện ở tầng staging của dbt (\texttt{stg\_ga4\_events}, \texttt{stg\_ga4\_items}). Nhờ tầng xử lý giữ nguyên cấu trúc gốc, kho dữ liệu chạy được trên cả dữ liệu công khai lẫn dữ liệu đã qua pipeline mà không phải sửa mô hình.
    \end{itemize}

    \item \textbf{Tầng kho dữ liệu (Load):} Dữ liệu JSONL được nạp vào BigQuery Sandbox, sau đó dbt dựng staging và mô hình sao gồm 4 bảng fact dùng chung 8 bảng dimension. Ba view staging làm phẳng tham số và mặt hàng, tách phiên và chọn giao dịch hợp lệ. Bốn bảng marts là \texttt{fct\_daily\_ecommerce}, \texttt{fct\_event\_counts}, \texttt{fct\_product\_performance} và \texttt{fct\_product\_channel\_daily}; chúng cung cấp các số đo phục vụ dashboard và ba KPI view. Sandbox được giữ trong giới hạn miễn phí và không yêu cầu bucket trung gian hoặc streaming inserts.

    \item \textbf{Tầng phân tích và ứng dụng:} Xây dựng dashboard trực quan, khai thác dữ liệu từ các bảng tổng hợp trong BigQuery. Dashboard cung cấp các chỉ số như số lượng sự kiện, số người dùng, số phiên, giao dịch, doanh thu và hiệu suất sản phẩm. Người dùng có thể theo dõi xu hướng theo thời gian và so sánh kết quả giữa các nhóm thiết bị, khu vực địa lý hoặc nguồn truy cập. Những thông tin này tạo cơ sở cho việc đánh giá hoạt động và hỗ trợ ra quyết định dựa trên dữ liệu.
\end{enumerate}

Luồng dữ liệu:
\[
\text{Nguồn}
\rightarrow \text{Trích xuất}
\rightarrow \text{Biến đổi và kiểm tra chất lượng}
\rightarrow \text{Kho dữ liệu}
\rightarrow \text{Dashboard/DSS}.
\]

\subsection{Hệ thống chỉ số (KPI)}

Mỗi chỉ số được định nghĩa một lần, ghi rõ công thức, đơn vị và mức tổng hợp, để mọi báo cáo dùng chung.

\subsection*{Nhóm A: Lưu lượng và mức độ tương tác}

\begin{table}[H]
    \centering
    \renewcommand{\arraystretch}{1.2}
    \caption{Các chỉ số về lưu lượng và mức độ tương tác}
    \label{tab:traffic_kpi}
    \begin{tabularx}{\textwidth}
    {
        >{\raggedright\arraybackslash}p{4.2cm}
        >{\raggedright\arraybackslash}X
    }
        \hline
        \textbf{Chỉ số} & \textbf{Định nghĩa} \\
        \hline
        Số người dùng &
        Số \texttt{user\_pseudo\_id} phân biệt. \\

        Số phiên &
        Số cặp (\texttt{user\_pseudo\_id}, \texttt{ga\_session\_id}) phân biệt. \\

        Số phiên/người dùng &
        Số phiên / số người dùng. \\

        Số lượt xem trang &
        Số sự kiện \texttt{page\_view}. \\

        Tỷ lệ phiên có tương tác &
        Số phiên có tương tác (theo quy tắc GA4, ví dụ \texttt{session\_engaged}) / tổng số phiên. \\

        Tỷ lệ người dùng mới &
        Số người dùng có sự kiện \texttt{first\_visit} / số người dùng. \\
        \hline
    \end{tabularx}
\end{table}

\subsection*{Nhóm B: Phễu chuyển đổi và hành vi mua hàng}

\begin{table}[H]
    \centering
    \renewcommand{\arraystretch}{1.2}
    \caption{Các chỉ số về phễu chuyển đổi và hành vi mua hàng}
    \label{tab:funnel_kpi}
    \begin{tabularx}{\textwidth}{
        >{\raggedright\arraybackslash}p{4.2cm}
        >{\raggedright\arraybackslash}X
    }
        \hline
        \textbf{Chỉ số} & \textbf{Định nghĩa} \\
        \hline
        Số lượt xem sản phẩm &
        Số sự kiện \texttt{view\_item}. \\

        Tỷ lệ thêm vào giỏ &
        Số phiên có \texttt{add\_to\_cart} / số phiên có \texttt{view\_item}. \\

        Tỷ lệ bắt đầu thanh toán &
        Số phiên có \texttt{begin\_checkout} / số phiên có \texttt{add\_to\_cart}. \\

        Tỷ lệ điền thông tin giao hàng &
        Số phiên có \texttt{add\_shipping\_info} / số phiên có \texttt{begin\_checkout}. \\

        Tỷ lệ điền thông tin thanh toán &
        Số phiên có \texttt{add\_payment\_info} / số phiên có \texttt{add\_shipping\_info}. \\

        Tỷ lệ hoàn tất mua hàng &
        Số phiên có giao dịch hợp lệ / số phiên có \texttt{add\_payment\_info}. \\

        Tỷ lệ rời bỏ giỏ hàng &
        $1 - \dfrac{\text{số phiên có giao dịch hợp lệ}}
        {\text{số phiên có \texttt{add\_to\_cart}}}$. \\

        Tỷ lệ chuyển đổi theo từng bước phễu &
        Số phiên đạt bước $k+1$ / số phiên đạt bước $k$. \\
        \hline
    \end{tabularx}
\end{table}

Phễu chuyển đổi gồm sáu bước: \texttt{view\_item} $\rightarrow$ \texttt{add\_to\_cart} $\rightarrow$ \texttt{begin\_checkout} $\rightarrow$ \texttt{add\_shipping\_info} $\rightarrow$ \texttt{add\_payment\_info} $\rightarrow$ mua hàng. Khảo sát dữ liệu xác nhận cả sáu bước đều có mặt trong 92 ngày; riêng \texttt{add\_to\_cart} vắng mặt trong 18/92 ngày nên biểu đồ theo ngày không nội suy cho các ngày thiếu và đánh dấu bằng cờ \texttt{is\_add\_to\_cart\_tracked} trong \texttt{dim\_date}.

\subsection*{Nhóm C: Chuyển đổi và doanh thu}

\begin{table}[H]
    \centering
    \renewcommand{\arraystretch}{1.2}
    \caption{Các chỉ số về chuyển đổi và doanh thu}
    \label{tab:revenue_kpi}
    \begin{tabularx}{\textwidth}{
        >{\raggedright\arraybackslash}p{4.2cm}
        >{\raggedright\arraybackslash}X
    }
        \hline
        \textbf{Chỉ số} & \textbf{Định nghĩa} \\
        \hline
        Số đơn hàng &
        Số \texttt{transaction\_id} phân biệt trong các giao dịch hợp lệ. \\

        Tỷ lệ chuyển đổi &
        Số phiên có ít nhất một giao dịch hợp lệ / tổng số phiên. \\

        Doanh thu &
        Tổng \texttt{purchase\_revenue\_in\_usd} của các giao dịch hợp lệ. \\

        Giá trị đơn trung bình (AOV) &
        Doanh thu / số đơn hàng. \\

        Doanh thu trên mỗi phiên &
        Doanh thu / số phiên. \\

        Doanh thu trên mỗi người dùng &
        Doanh thu / số người dùng. \\

        Tỷ trọng doanh thu theo kênh &
        Doanh thu của kênh / tổng doanh thu. \\
        \hline
    \end{tabularx}
\end{table}

\subsection*{Định nghĩa giao dịch hợp lệ}

Các chỉ số mua hàng ở nhóm B và C chỉ tính trên \textbf{giao dịch hợp lệ}, được định nghĩa bởi đồng thời bốn điều kiện:

\begin{itemize}
    \item Là sự kiện \texttt{purchase}.
    \item Có \texttt{transaction\_id} khác \texttt{NULL} và không thuộc các giá trị giữ chỗ \texttt{'(not set)'}, \texttt{'<Other>'}, \texttt{'(data deleted)'} hoặc chuỗi rỗng.
    \item Có doanh thu \texttt{purchase\_revenue\_in\_usd} lớn hơn 0.
    \item Nếu một người dùng (\texttt{user\_pseudo\_id}) có phát sinh được nhiều sự kiện có trùng \texttt{transaction\_id}, chỉ giữ lấy sự kiện có thuộc tính \texttt{event\_timestamp} sớm nhất.
\end{itemize}

Định nghĩa này được cài đặt một lần duy nhất ở mô hình staging \texttt{stg\_valid\_orders} và được mọi bảng fact, bảng tổng hợp và KPI kế thừa; các bản ghi bị loại theo từng điều kiện được ghi vào nhật ký chất lượng dữ liệu. \medskip

Chuyển đổi là sự kiện \texttt{purchase} được ghi nhận theo nguồn của phiên hoặc theo nguồn của người dùng nếu không xác định được nguồn phiên. Quy tắc này giải quyết rủi ro so sánh không đồng nhất nêu ở mục 2.1 và 2.2. Nhóm sẽ ghi rõ mô hình phân bổ được sử dụng.

\subsection{Chức năng hỗ trợ ra quyết định}

Ứng dụng dashboard gồm các chức năng sau:

\begin{enumerate}
    \item \textbf{Tổng quan điều hành:} Các KPI chính, xu hướng theo ngày/tuần/tháng, so sánh với kỳ trước.

    \item \textbf{Bộ lọc đa chiều:} Theo khoảng thời gian, kênh, nguồn/medium, chiến dịch, thiết bị, quốc gia/khu vực, nhóm sản phẩm và loại người dùng. Bộ lọc nhóm sản phẩm chỉ tác động đến dữ liệu sản phẩm; quốc gia và loại người dùng chỉ áp dụng cho các số liệu theo phiên nơi có các chiều tương ứng.

    \item \textbf{So sánh hiệu suất:} Đặt nhiều kênh, chiến dịch hoặc kỳ cạnh nhau trên cùng bộ chỉ số; xếp hạng theo doanh thu, tỷ lệ chuyển đổi và doanh thu trên mỗi phiên.

    \item \textbf{Phân tích chất lượng kênh:} Đối chiếu quy mô lưu lượng với chất lượng lưu lượng (tương tác, chuyển đổi, doanh thu/phiên), giúp phân biệt kênh ``nhiều truy cập nhưng ít đơn'' với kênh ``ít truy cập nhưng chuyển đổi tốt''.

    \item \textbf{Phân tích phễu chuyển đổi:} Từ xem sản phẩm đến mua hàng, chỉ ra bước có tỷ lệ rời bỏ cao nhất; có thể chia theo kênh, thiết bị hoặc nhóm sản phẩm.

    \item \textbf{Phân tích chi tiết:} Từ kênh xuống chiến dịch, từ tháng xuống ngày, từ nhóm sản phẩm xuống sản phẩm.

    \item \textbf{Phân khúc khách hàng:} Nhóm theo thiết bị, vùng, nguồn truy cập và hành vi (người dùng mới/quay lại, mức độ tương tác).

    \item \textbf{Cảnh báo theo ngưỡng:} Đánh dấu kênh/chiến dịch có chỉ số bất thường (tỷ lệ chuyển đổi hoặc doanh thu, phiên giảm đáng kể so với trung bình kỳ trước, tỷ lệ rời bỏ giỏ hàng tăng, lưu lượng tăng đột biến nhưng chuyển đổi không tăng theo). Ngưỡng do người dùng cấu hình.

    \item \textbf{Mô phỏng kịch bản:} Người dùng điều chỉnh giả định về lưu lượng của từng kênh hoặc tỷ lệ chuyển đổi rồi xem ước tính ảnh hưởng đến số đơn và doanh thu dựa trên hiệu suất lịch sử.

    \item \textbf{Truy vết số liệu:} Từ một con số trên dashboard xem được công thức và các bản ghi nguồn tạo ra nó.
\end{enumerate}

\subsection{Ánh xạ chức năng giao diện với người ra quyết định}

Bảng dưới đây phản ánh cách giao diện nhóm các trang theo vai trò. Mục đích là giúp người đọc nhận ra ngay phần nào của ứng dụng được thiết kế cho quyết định điều hành, theo dõi chiến dịch hay phân tích chuyên sâu về hành vi mua hàng. Bộ lọc đa chiều là thành phần dùng chung trên các trang, nên hệ thống có mười chức năng nhưng chín trang giao diện. \medskip

\begin{table}[H]
    \centering
    \renewcommand{\arraystretch}{1.25}
    \caption{Chức năng giao diện theo từng nhóm người ra quyết định}
    \label{tab:dashboard-role-mapping}
    \begin{tabularx}{\textwidth}{
        >{\raggedright\arraybackslash}p{3.8cm}
        >{\raggedright\arraybackslash}p{2.0cm}
        >{\raggedright\arraybackslash}p{3.7cm}
        X
    }
        \hline
        \textbf{Chức năng} &
        \textbf{CMO} &
        \textbf{Digital Marketing Manager} &
        \textbf{E-commerce \& Content Specialist} \\
        \hline
        Tổng quan điều hành & Có & - & - \\
        Bộ lọc đa chiều & Dùng chung & Dùng chung & Dùng chung \\
        So sánh hiệu suất & Có & Có & - \\
        Phân tích chất lượng kênh & Có & Có & - \\
        Phân tích phễu chuyển đổi & - & - & Có \\
        Phân tích chi tiết & - & Có & Có \\
        Phân khúc khách hàng & - & - & Có \\
        Cảnh báo theo ngưỡng & Có & Có & - \\
        Mô phỏng kịch bản & Có & Có & - \\
        Truy vết số liệu & Có & Có & Có \\
        \hline
    \end{tabularx}
\end{table}

Chế độ ``CMO'' mở đầu tiên trang Tổng quan, ``Manager'' mở trang So sánh hiệu suất, còn ``Specialist'' mở trang Phễu chuyển đổi. Chế độ ``Tất cả'' hiển thị đủ chín trang để người dùng xem ngoài nhóm chức năng mặc định của vai trò mình, chẳng hạn CMO cần xem phân khúc khách hàng. Các vai trò chỉ thay đổi điều hướng và nhóm trang được trình bày, không phải cơ chế xác thực hoặc phân quyền trên dữ liệu. \medskip

\subsection{Ánh xạ giải pháp với vấn đề và người ra quyết định}

\begin{table}[H]
    \centering
    \renewcommand{\arraystretch}{1.3}
    \caption{Ánh xạ giữa vấn đề kinh doanh, giải pháp và người ra quyết định}
    \label{tab:solution_mapping}
    \begin{tabularx}{\textwidth}{
        >{\raggedright\arraybackslash}p{2.5cm}
        >{\raggedright\arraybackslash}X
        >{\centering\arraybackslash}p{1.8cm}
        >{\raggedright\arraybackslash}p{3cm}
    }
        \hline
        \textbf{Vấn đề} &
        \textbf{Thành phần giải pháp} &
        \textbf{Mức độ} &
        \textbf{Người dùng chính} \\
        \hline

        1. Dữ liệu phân tán &
        Kho dữ liệu tập trung và ETL chuẩn hóa &
        Trọng tâm &
        Tất cả \\

        2. Khó đánh giá hiệu quả chiến dịch &
        Hệ KPI thống nhất; so sánh hiệu suất; phân tích chất lượng kênh &
        Trọng tâm &
        CMO, Digital Marketing Manager \\

        3. Khó phân tích hành vi, phễu &
        Phân tích phễu theo phiên; drill-down &
        Trọng tâm &
        E-commerce \& Content Specialist \\

        4. Khó phân khúc khách hàng &
        Phân tích theo thiết bị, vùng, nguồn, hành vi &
        Một phần &
        CMO, E-commerce \& Content Specialist \\

        5. Báo cáo chậm, thiếu nhất quán &
        ETL tự động; dashboard có bộ lọc; định nghĩa KPI duy nhất &
        Trọng tâm &
        Digital Marketing Manager \\

        6. Khó dự báo và phân bổ ngân sách &
        Phân tích chất lượng kênh làm cơ sở ưu tiên nguồn lực; mô phỏng kịch bản theo lưu lượng và chuyển đổi &
        Một phần &
        CMO \\

        \hline
    \end{tabularx}
\end{table}


Ánh xạ này cho thấy loại hỗ trợ mà mỗi tầng cung cấp cho người ra quyết định. Các chỉ số được tính từ cùng một định nghĩa giao dịch và cùng các dimension dùng chung, thay vì từ các bảng tùy ý không đồng bộ. Mức độ tin cậy của kết quả lượng hóa được trình bày trong phần Evaluation: quá trình thử hiện tại sử dụng một mini-batch Sandbox, không phải đánh giá hiệu năng đã hoàn tất trên cả ba tháng dữ liệu.

\newpage
\subsection{Phạm vi, giả định và giới hạn}

\subsection*{Phạm vi.}
Hệ thống phát lại một khoảng dữ liệu lịch sử có kiểm soát để thử nghiệm luồng tương tự tiếp nhận trực tuyến. Publisher sắp xếp các sự kiện theo thời gian trước khi đưa qua emulator, DirectRunner làm sạch và ghi JSONL, rồi Sandbox lưu raw và dbt tổng hợp. Khối lượng đánh giá mô phỏng được giới hạn để không vượt quota. Các hướng mở rộng lên hạ tầng quản lý hoặc nguồn dữ liệu trả phí được tách riêng trong phần kết luận, không nằm trong luồng mặc định không bật billing.\\

\subsection*{Giả định.}
\begin{itemize}
    \item Một phiên được xác định bởi cặp \texttt{user\_pseudo\_id} và \texttt{ga\_session\_id}; kênh của phiên được suy ra từ source/medium theo bộ quy tắc gom nhóm do nhóm định nghĩa và công bố.

    \item Doanh thu lấy từ giá trị giao dịch trong dữ liệu; chưa có dữ liệu giá vốn nên chưa tính được lợi nhuận. Vì vậy, doanh thu cao chưa có nghĩa là lợi nhuận cao, đúng như lưu ý ở mục 1.1.
\end{itemize}

\subsection*{Giới hạn.}
\begin{itemize}
    \item Không có dữ liệu chi phí, hiển thị và lượt nhấp từ nền tảng quảng cáo nên không đánh giá được CPA, ROAS, CPC và hiệu quả đầu tư theo chi phí.

    \item Do dữ liệu bị làm mờ, một số kết quả (như phân bổ nguồn, quy mô doanh thu) có thể lệch so với thực tế; nhóm ghi nhận tỷ lệ bản ghi bị ảnh hưởng.

    \item Kết quả mô phỏng kịch bản dựa trên hiệu suất lịch sử và giả định quan hệ ổn định giữa lưu lượng, chuyển đổi và doanh thu, nên chỉ mang tính tham khảo.

    \item Lợi ích kỳ vọng chưa được đo trên quy trình thực; nhóm đánh giá hệ thống qua kiểm thử tính đúng đắn của dữ liệu, tính nhất quán KPI, hiệu năng truy vấn và mức độ đáp ứng các câu hỏi nghiệp vụ.
\end{itemize}
